
\setlength{\parindent}{10pt}

\chapter*[Preface]{Preface}

\emph{\color{blue}Elementary Differential Equations with Boundary Value Problems}
 is
written for students in science, engineering, and mathematics who have
completed calculus  through partial  differentiation. If  your syllabus
includes \zcref{ch:systems} (\nameref*{ch:systems}),
your students should have some preparation in linear algebra.

In writing this book I have been guided by the these principles:

\begin{itemize}

\item An elementary text should be written so  the student can
read it with comprehension without too much pain. I have tried to
put myself in the student's place, and have chosen to err on the side
of too much detail rather than not enough.

\item An elementary text can't be better than its exercises.
This text includes  1695 numbered exercises, many with several parts.
They range in difficulty from routine to very challenging.

\item An elementary text should be written in an informal but
mathematically accurate way, illustrated by appropriate graphics.
I have tried to formulate
mathematical concepts succinctly
 in language that students can understand. I have
minimized the number of explicitly stated theorems and definitions,
preferring to deal with concepts in a more conversational way,
copiously illustrated by 250 completely worked out examples. Where
appropriate, concepts and results are depicted in 144 figures.
\end{itemize}

Although I believe that the computer is an immensely valuable tool for
learning, doing, and writing mathematics, the selection and treatment
of topics in this text reflects my pedagogical orientation along
traditional lines. However, I have incorporated what I believe to be
the best use of modern technology, so you can select the level of
technology that you want to include in your course. The text includes
336 exercises --- identified by the symbols \Cex  and \CGex --- that call
for graphics or computation and graphics. There are also
73 laboratory exercises --- identified by \Lex --- that require extensive
use of technology. In addition, several sections include informal advice on
the use of technology. If you prefer not to emphasize technology,
simply ignore these exercises and the  advice.

There are two schools of thought on whether techniques and
applications should be treated together or separately. I have chosen
to separate them;
thus, \zcref{ch:first_order} deals with techniques for solving
first order equations, and \zcref{ch:apps_first} deals with applications.
Similarly,
\zcref{ch:second_order} deals with techniques for solving second order
equations, and \zcref{ch:apps_second} deals with applications. However, the
exercise sets of the sections dealing with techniques  include some
applied problems.

Traditionally oriented elementary differential
equations texts are occasionally  criticized as being collections of
unrelated methods for solving miscellaneous problems. To some extent
this is true;     after all, no single method  applies to all
situations. Nevertheless, I believe that one idea can go a long way
toward unifying some of the techniques for solving diverse problems:
variation of parameters. I use variation of parameters at the earliest
opportunity in \zcref{sec:linear_first_order},
to solve the nonhomogeneous linear equation,
given a nontrivial solution of the complementary equation.
You may find this annoying, since most of us learned that one should
use integrating factors for this task, while perhaps mentioning the
variation of parameters option in an exercise. However, there's
little difference between the two approaches, since an integrating
factor is nothing more than the reciprocal of a nontrivial solution of
the complementary equation. The advantage of using variation of
parameters here is that it introduces the concept in its simplest form
and focuses the student's attention on the idea of seeking a solution
$y$ of a differential equation by writing it as $y=uy_1$, where $y_1$
is a known solution of related equation and $u$ is a function to be
determined. I use this idea in nonstandard ways, as follows:

\begin{itemize}
\item
In \zcref{sec:transform_separable} to solve nonlinear first order equations,
such as Bernoulli equations and nonlinear homogeneous equations.

\item  In \zcref{ch:numerical}  for  numerical solution of
semilinear  first order equations.

\item In \zcref{sec:const_coeff_homog} to avoid the necessity of introducing
complex
exponentials in solving a second order constant coefficient
homogeneous equation with characteristic polynomials  that have complex
zeros.

\item In \zcref{sec:undetermined_coeffs_i,sec:undetermined_coeffs_ii,sec:higher_undetermined}
for the method of undetermined
coefficients. (If the method of annihilators is your preferred
approach to this problem, compare the labor involved in solving, for
example, $y''+y'+y=x^4e^x$ by the method of annihilators and the
method used in \zcref{sec:undetermined_coeffs_i}.)
\end{itemize}

Introducing variation of parameters as early as possible
(\zcref{sec:linear_first_order})
prepares the student for the concept when it appears again in more
complex forms in \zcref{sec:reduction_order}, where reduction of order is used
not merely to find a second solution of the complementary equation, but
also to find the general solution of the nonhomogeneous equation, and in 
\zcref{sec:variation_param,sec:higher_variation_param,sec:systems_variation_param},
that treat the usual variation of parameters problem for second and higher
order linear equations and for linear systems.

\iftoggle{boundaryvalues}{%
\zcref[S]{ch:bvp_fourier} develops the theory of Fourier series.
\zcref[S]{sec:eig_probs} discusses the five main eigenvalue problems
that arise in connection with the method of separation of
variables for the heat and wave equations and for Laplace's
equation over a rectangular domain:
\begin{enumext}[label=\arabic*,wrap-label={Problem #1:},labelwidth=5em]
\item $y''+\lambda y=0,\quad y(0)=0,\quad y(L)=0$
\item $y''+\lambda y=0,\quad y'(0)=0,\quad y'(L)=0$
\item $y''+\lambda y=0,\quad y(0)=0,\quad y'(L)=0$
\item $y''+\lambda y=0,\quad y'(0)=0,\quad y(L)=0$
\item $y''+\lambda y=0,\quad y(-L)=y(L),\quad y'(-L)=y'(L)$
\end{enumext}

These problems are handled in a
unified way     for example, a single theorem shows that the
eigenvalues of all five problems are nonnegative.

\zcref[S]{sec:fourier_i} presents the  Fourier series expansion
of functions defined on
on $[-L,L]$,   interpreting it as an expansion in terms of the
eigenfunctions of Problem~5.

\zcref[S]{sec:fourier_ii}  presents the
 Fourier sine and cosine expansions of functions defined on
$[0,L]$, interpreting them as expansions in terms of the
eigenfunctions of Problems~1 and 2, respectively. In addition,
\zcref{sec:fourier_i} includes what I call the mixed Fourier sine and
cosine expansions, in terms of the eigenfunctions of Problems~4
and 5, respectively. In all cases, the convergence properties of
these series are deduced from the convergence properties of
the Fourier series discussed in \zcref[S]{sec:eig_probs}.

\zcref[S]{ch:fourier_pde} consists of four sections devoted to the
heat equation, the wave equation, and Laplace's equation
in rectangular and polar coordinates.
For all three, I consider homogeneous boundary conditions
of the four types occurring in Problems~1--4. I present
the method of separation of variables as a way of
choosing the appropriate form for the series expansion
of the solution of the given problem, stating---without
belaboring the point---that the expansion may fall short of being
an actual solution, and giving an indication of conditions under
which the formal solution is an actual solution. In particular, I
found it necessary
to devote some detail to this question  in connection with the wave
equation in \zcref{sec:wave_eq}.

In Sections \zcref[noname]{sec:heat_eq} (\nameref*{sec:heat_eq}) and
\zcref[noname]{sec:wave_eq} (\nameref*{sec:wave_eq}) I
%In \zcref{sec:heat_eq,sec:wave_eq} (\nameref*{sec:heat_eq} and \nameref*{sec:wave_eq}, respectively) I
devote considerable effort to devising examples and numerous
exercises where the functions defining the initial conditions
satisfy
the homogeneous boundary conditions. Similarly, in most of the
examples and exercises \zcref{sec:laplace_rect}
(\nameref*{sec:laplace_rect}), the functions defining the boundary conditions
on a given side of the rectangular domain satisfy homogeneous
boundary conditions at the endpoints of the same type (Dirichlet or
Neumann) as the boundary conditions imposed on adjacent sides of
the region. Therefore  the formal solutions obtained in many of the
examples and exercises are  actual solutions.

\zcref[S]{sec:twopoint_bvp} deals with two-point value problems for
a second order ordinary differential equation. Conditions for
existence
and uniqueness of solutions are given, and the construction
of Green's functions is included.

\zcref[S]{sec:sturm_liouville} presents the elementary aspects of Sturm-Liouville
theory.
}{}% end toggle{boundaryvalues}

You may also find the following to be of interest:

\begin{itemize}

\item \zcref[S]{sec:integrating_factors} deals with integrating factors of the
form $\mu=p(x)q(y)$, in addition to those of the form $\mu=p(x)$ and
$\mu=q(y)$ discussed in most texts.

\item \zcref[S]{sec:auto_second_order} makes phase plane analysis of nonlinear
second order autonomous equations accessible to students who have not taken
linear algebra, since eigenvalues and eigenvectors do not enter into the
treatment. Phase plane analysis of constant coefficient linear systems
is included in
\zcref[range]{sec:system_const_coeff_i,sec:system_const_coeff_iii}.

\item \zcref[S]{sec:app_curves} presents an extensive discussion of
applications of differential equations to curves.

\item \zcref[S]{sec:central_force} studies motion under a central force,
which may be useful to students interested in the mathematics of
satellite orbits.

\item \zcref[S,range]{sec:frobenius_i,sec:frobenius_iii} present the method of
Frobenius in more detail than in most texts. The approach is to systematize the
computations in a way that avoids the necessity of substituting the
unknown Frobenius series into each equation. This leads to efficiency
in the computation of the coefficients of the Frobenius solution. It
also clarifies the case where the roots of the indicial equation
differ by an integer (\zcref{sec:frobenius_iii}).

\item The free Student Solutions Manual contains solutions
of most of the even-numbered exercises.

%\item The free Instructor's Solutions Manual is available by email to
%\href{mailto:wtrench@trinity.edu}{wtrench@trinity.edu},
%subject to verification of the requestor's faculty status.

\end{itemize}

The following observations may be helpful as you choose your
 syllabus:

\begin{itemize}
\item \zcref[S]{sec:exist_unique}  is the only specific prerequisite for
\zcref{ch:numerical}. To accommodate institutions that offer a separate course
in numerical analysis,  \zcref*{ch:numerical} is not
a prerequisite for any other section in the text.

\item The sections in \zcref{ch:apps_first} are independent of each
other, and are not prerequisites for any of the later chapters. This
is also true of the sections in \zcref{ch:apps_second}, except that
\zcref{sec:spring_i} is a prerequisite for \zcref{sec:spring_ii}.

\item \zcref[S]{ch:series,ch:laplace,ch:higher_order} can be covered in any
order after the topics selected  from \zcref{ch:second_order}. For example, you
can
proceed directly from \zcref{ch:second_order} to \zcref{ch:higher_order}.

\item The second order Euler equation is discussed in
\zcref{sec:regular_singular}, where it sets the stage for the method of
Frobenius. As noted at the beginning of \zcref*{sec:regular_singular}, if you
want to include Euler equations in
your syllabus while omitting the method of Frobenius, you can skip the
introductory paragraphs in \zcref*{sec:regular_singular} and begin with
\zcref{thmtype:7.4.2}. You can then cover
\zcref*{sec:regular_singular} immediately after \zcref{sec:const_coeff_homog}.

\iftoggle{boundaryvalues}{%
\item \zcref[S,nocomp]{ch:bvp_fourier,ch:fourier_pde,ch:bvp_second_order} can be covered at any time
after the completion of \zcref{ch:second_order}.%
}{}% end toggle{boundaryvalues}

\end{itemize}
\medskip

\hfill William F. Trench\hspace{1in}
